% Generated draft: requires mathematical and source review.
\documentclass[11pt]{article}
\usepackage[margin=25mm]{geometry}
\usepackage{fontspec}
\setmainfont{Latin Modern Roman}
\usepackage{amsmath,amssymb,mathtools,unicode-math}
\setmathfont{Latin Modern Math}
\usepackage{enumitem,needspace,xurl,hyperref}
\setlength{\parindent}{0pt}
\setlength{\parskip}{5pt}
\setlength{\emergencystretch}{3em}
\newcommand{\N}{\mathbb{N}}
\newcommand{\Z}{\mathbb{Z}}
\newcommand{\Q}{\mathbb{Q}}
\newcommand{\R}{\mathbb{R}}
\newcommand{\C}{\mathbb{C}}
\newcommand{\F}{\mathbb{F}}
\newcommand{\A}{\mathbb{A}}
\newcommand{\PP}{\mathbb{P}}
\newcommand{\E}{\mathbb{E}}
\newcommand{\eps}{\varepsilon}
\newcommand{\dd}{\,\mathrm{d}}
\begin{document}
\section*{UnsolvedMath ID: 800668}
\textbf{Source problem code:} catalogue entry
\subsection{Definitions and mathematical statement}
Under the conventional real-valued interpretation, for a topological space \(X\), let \(C_p(X)=C(X,\R)\) have the topology of pointwise convergence, and let \(\mathbf{0}\) denote the constant zero function. Thus, for \(A\subseteq C_p(X)\),
\[
\mathbf{0}\in\overline A
\quad\Longleftrightarrow\quad
\forall F\subseteq X\text{ finite}\ \forall\varepsilon>0\ \exists f\in A\ \forall x\in F,\ |f(x)|<\varepsilon.
\]
Use \(\mathbb N=\{1,2,\ldots\}\).

In the standard definition, \(C_p(X)\) has countable fan tightness at \(\mathbf{0}\) if, for every sequence \((A_n)_{n\in\mathbb N}\) of subsets of \(C_p(X)\) satisfying \(\mathbf{0}\in\overline{A_n}\) for every \(n\), there are finite sets \(B_n\subseteq A_n\) such that
\[
\mathbf{0}\in\overline{\bigcup_{n\in\mathbb N}B_n}.
\]
Translation by any continuous function is a homeomorphism of \(C_p(X)\), so this condition at \(\mathbf{0}\) is equivalent to the corresponding condition at every point.

The supplied excerpt defines the weak sequence selection property by modifying the sequence selection property (SSP), but does not include the definition of SSP. Under the conventional reading that the input to SSP consists of sequences converging to \(\mathbf{0}\), its weak version is the following property: for every array \((f^n_m)_{n,m\in\mathbb N}\) in \(C_p(X)\) satisfying
\[
\forall n\in\mathbb N\ \forall x\in X,\qquad
\lim_{m\to\infty}f^n_m(x)=0,
\]
there is a sequence \((k_n)_{n\in\mathbb N}\in\mathbb N^{\mathbb N}\) such that
\[
\mathbf{0}\in\overline{\{f^n_{k_n}:n\in\mathbb N\}}.
\]
Equivalently, the chosen indices must satisfy
\[
\forall F\subseteq X\text{ finite}\ \forall\varepsilon>0\ \exists n\in\mathbb N\ \forall x\in F,
\qquad |f^n_{k_n}(x)|<\varepsilon.
\]

\textbf{Question (supplied Problem 41).} Does countable fan tightness of \(C_p(X)\) imply this weak sequence selection property?

\textbf{Unresolved scope qualifications.} The excerpt does not specify the admissible class of spaces \(X\), the scalar field in the definition of \(C_p(X)\), or the full SSP definition. The definitions above explicitly adopt the conventional real-valued interpretation and the stated interpretation of SSP; they require confirmation against the source context. No separation axiom for \(X\) is imposed here. The intended universal quantifier over \(X\) remains dependent on the source's missing ambient hypotheses.
\subsection{Short English statement}
The question asks whether countable fan tightness of \(C_p(X)\) guarantees that one can choose one function from each of countably many sequences converging pointwise to zero so that zero lies in the closure of the chosen functions. The excerpt leaves the permitted spaces \(X\) and the full SSP definition unspecified, so the displayed formulation is conditional on the conventions stated above.
\subsection{Sources}
\begin{itemize}
\item Local UnsolvedMath registry, permanent ID \(800668\), exact source problem code \texttt{catalogue entry}; title: \emph{Selection Principles Milestone Problems --- Problem 41 ([CpHure, Problem 1])}. The registry lists Tsaban as proposer and \(2003\) as the proposed year. Its supplied status is \texttt{open}; this is a registry label, not an independently verified current mathematical assessment.
\item Tsaban, \emph{Selection Principles in Mathematics: A Milestone of Open Problems} (2003), Problem 41. Supplied source URL: \url{https://arxiv.org/abs/math/0312182}. The record reports access on 2026-10-16 and extraction by \texttt{arxiv\_latex}; these are supplied provenance details, not a claim of source access or verification for this draft.
\item The supplied excerpt attributes the weak sequence selection property to the reference label [CpHure], and its background identifies [CpHure, Problem 1] as attribution/context. No full bibliographic entry for [CpHure] is supplied.
\end{itemize}
\end{document}
